\documentclass[11pt,a4paper]{article}
\usepackage[T1]{fontenc}
\usepackage[utf8]{inputenc}
\usepackage{lmodern}
\usepackage{amsmath,amssymb}
\usepackage{geometry}
\usepackage{microtype}
\usepackage{enumitem}
\usepackage{parskip}
\usepackage{titlesec}

\geometry{top=2.5cm, bottom=2.5cm, left=2.8cm, right=2.8cm}
\titlespacing*{\section}{0pt}{1.6ex}{0.6ex}
\setlength{\parskip}{5pt}

\begin{document}

\begin{center}
{\large\textbf{EE 533 -- Assignment 2}} \quad Kerem Recep Gür \quad 2448462
\end{center}
\medskip\hrule\bigskip

%-----------------------------------------------------------------------
\section*{Problem 1}

\[
  p_X(x)=\tau(1-\tau)^x\ (x=0,1,2,\ldots),
  \quad
  p_Y(y)=s(1-s)^{y-1}\ (y=1,2,3,\ldots),
  \quad
  p_Z(z)=\tfrac{1}{k}\ (z=1,\ldots,k).
\]

\medskip
\textbf{(a) $D(p_X\|p_Y) = +\infty$.}

$p_X$ is defined at $x=0$ (its value is $\tau > 0$),
but $p_Y$ is \emph{not} defined at $x=0$ (its domain starts at 1).
We treat $p_Y(0) = 0$.
The $x=0$ term inside the KL sum then becomes:
\[
  p_X(0)\cdot\ln\frac{p_X(0)}{p_Y(0)}
  = \tau\cdot\ln\frac{\tau}{0}
  = \tau\cdot(+\infty)
  = +\infty.
\]
Since one term is already $+\infty$, the whole sum is $+\infty$.

\medskip
\textbf{(b) $D(p_Y\|p_X)$.}

Substitute the PMFs and expand the log:
\begin{align*}
  D(p_Y\|p_X)
  &= \sum_{y=1}^{\infty} s(1-s)^{y-1}
     \bigl[\ln s - \ln\tau + (y-1)\ln(1-s) - y\ln(1-\tau)\bigr]\\
  &= \ln\!\tfrac{s}{\tau}\underbrace{\sum p_Y}_{1}
   + \ln(1-s)\underbrace{\mathbb{E}[Y-1]}_{(1-s)/s}
   - \ln(1-\tau)\underbrace{\mathbb{E}[Y]}_{1/s}.
\end{align*}
\[
  \boxed{D(p_Y\|p_X) = \ln\frac{s}{\tau} + \frac{1-s}{s}\ln(1-s) - \frac{1}{s}\ln(1-\tau).}
\]

\medskip
\textbf{(c) $D(p_Z\|p_Y)$.}

\begin{align*}
  D(p_Z\|p_Y)
  &= \frac{1}{k}\sum_{z=1}^{k}\ln\frac{1/k}{s(1-s)^{z-1}}
   = \frac{1}{k}\sum_{z=1}^{k}\bigl[-\ln k-\ln s-(z-1)\ln(1-s)\bigr].
\end{align*}
Using $\sum_{z=1}^{k}(z-1)=k(k-1)/2$:
\[
  \boxed{D(p_Z\|p_Y) = -\ln k - \ln s - \frac{k-1}{2}\ln(1-s).}
\]

%-----------------------------------------------------------------------
\section*{Problem 2}

We show $H(X_0\mid X_n)\le H(X_0\mid X_{n+1})$ for every $n$.

\textbf{Step 1.} The Markov property means $X_0\perp X_{n+1}\mid X_n$,
so knowing $X_{n+1}$ on top of $X_n$ adds nothing about $X_0$:
\[
  H(X_0\mid X_n,\;X_{n+1}) = H(X_0\mid X_n). \tag{1}
\]

\textbf{Step 2.} Conditioning on more variables never increases entropy:
\[
  H(X_0\mid X_n,\;X_{n+1}) \le H(X_0\mid X_{n+1}). \tag{2}
\]

Combining (1) and (2):
\[
  H(X_0\mid X_n)
  \;\overset{(1)}{=}\;
  H(X_0\mid X_n,X_{n+1})
  \;\overset{(2)}{\le}\;
  H(X_0\mid X_{n+1}).
  \qquad\square
\]

%-----------------------------------------------------------------------
\section*{Problem 3}

By the Law of Large Numbers, $\frac{1}{n}\sum_{i=1}^n \ln X_i \to \mathbb{E}[\ln X]$.
Since $(X_1\cdots X_n)^{1/n} = e^{\frac{1}{n}\sum \ln X_i}$, the limit is $e^{\mathbb{E}[\ln X]}$.

\[
  \mathbb{E}[\ln X]
  = \tfrac{1}{2}\ln 1 + \tfrac{1}{4}\ln 2 + \tfrac{1}{4}\ln 3
  = \tfrac{\ln 2 + \ln 3}{4}
  = \tfrac{\ln 6}{4}.
\]

\[
  \boxed{(X_1 X_2\cdots X_n)^{1/n} \;\xrightarrow{\;n\to\infty\;}\; e^{(\ln 6)/4} = \sqrt[4]{6} \approx 1.565.}
\]

%-----------------------------------------------------------------------
\section*{Problem 4}

\textbf{(a)} Split $2n$ samples into two halves; let $\hat{p}_n,\hat{p}_n'$ be their empirical PMFs.
Then $\hat{p}_{2n}=\frac{1}{2}\hat{p}_n+\frac{1}{2}\hat{p}_n'$.
By convexity of $D(\cdot\|p)$:
\[
  E[D(\hat{p}_{2n}\|p)]
  \le \tfrac{1}{2}E[D(\hat{p}_n\|p)]+\tfrac{1}{2}E[D(\hat{p}_n'\|p)].
\]
Since both halves are drawn i.i.d.\ from $p$, they are identically distributed,
so $E[D(\hat{p}_n'\|p)]=E[D(\hat{p}_n\|p)]$, giving
$E[D(\hat{p}_{2n}\|p)]\le E[D(\hat{p}_n\|p)]$.

\textbf{(b)} Let $\hat{p}_{n-1}^{(i)}$ be the empirical PMF of the $n-1$ samples obtained by deleting $X_i$.
Then $\hat{p}_n=\frac{1}{n}\sum_{i=1}^{n}\hat{p}_{n-1}^{(i)}$.
By convexity of $D(\cdot\|p)$:
\[
  E[D(\hat{p}_n\|p)]
  \le \frac{1}{n}\sum_{i=1}^{n}E[D(\hat{p}_{n-1}^{(i)}\|p)].
\]
Each $\hat{p}_{n-1}^{(i)}$ consists of $n-1$ i.i.d.\ samples from $p$,
so all terms are equal to $E[D(\hat{p}_{n-1}\|p)]$, giving
$E[D(\hat{p}_n\|p)]\le E[D(\hat{p}_{n-1}\|p)]$.

%-----------------------------------------------------------------------
\section*{Problem 5}

\textbf{$C_1=\{00,01,0\}$:} \texttt{0} is a prefix of \texttt{00}, so not instantaneous.
String \texttt{00} decodes as $\langle\texttt{00}\rangle$ or $\langle\texttt{0,0}\rangle$, so not uniquely decodable.

\textbf{$C_2=\{00,01,100,101,11\}$:} Length-2 words end with \texttt{00/01/11}, length-3 words start with \texttt{10} — no codeword is a prefix of another. Instantaneous $\Rightarrow$ uniquely decodable.

\textbf{$C_3=\{0,10,110,1110,\ldots\}$:} Each codeword ends in \texttt{0}; the next longer one has a \texttt{1} at that position. No codeword is a prefix of another. Instantaneous $\Rightarrow$ uniquely decodable.

\textbf{$C_4=\{0,00,000,0000\}$:} \texttt{0} is a prefix of all others, so not instantaneous.
String \texttt{0000} decodes as $\langle\texttt{0000}\rangle$ or $\langle\texttt{00,00}\rangle$ or $\langle\texttt{0,0,0,0}\rangle$, etc., so not uniquely decodable.

\bigskip
\begin{center}
\begin{tabular}{lcc}
\hline
Code & Uniquely decodable & Instantaneous \\
\hline
$C_1$ & No  & No  \\
$C_2$ & Yes & Yes \\
$C_3$ & Yes & Yes \\
$C_4$ & No  & No  \\
\hline
\end{tabular}
\end{center}

\end{document}
